% -*- root: ../main.tex -*-

% Esporre l'obiettivo del progetto dandone una visione complessiva. Devono essere illustrate le caratteristiche salienti del progetto; deve essere chiara la distinzione tra le tecnologie usate/assemblate durante lo svolgimento dell'elaborato e il contributo tecnologico/scientifico e effettivamente apportato dal gruppo.
% 3000 - 6000 battute

\chapter{Introduzione}
ISIQuiz è un progetto che nasce con l'intento di sviluppare un \textbf{gioco a quiz} che permetta ad uno studente di ripassare il contenuto dei corsi di \textbf{Ingegneria e Scienze Informatiche} attraverso domande a scelta multipla che coprono una o più materie di suo interesse. Lo scopo di \textbf{ISIQuiz} è quindi quello di gamificare il ripasso pre-esame attraverso uno strumento stimolante che permetta agli studenti di tracciare i propri progressi.
%img
\begin{figure}[H]
    \label{fig:Logo}
    \centering
    \includegraphics[width=0.8\textwidth]{Extra/ISIQuizLogoLineTransparent.png}
    \caption{Il logo del progetto}
\end{figure}

\section{Overview}
L'idea è nata dall'osservazione degli incontri tra studenti nelle giornate antecedenti ad una prova orale o scritta. Per preparasi al meglio sulle nozioni teoriche, nel gruppo si crea una dinamica simile a docente-studente, dove a turno ognuno fa domande di carattere generale su alcuni argomenti alle quali gli altri presenti cercano di rispondere. In questo modo, lo studente che risponde alle domande è in grado di ripassare e, con un numero sufficiente di quesiti, coprire gli argomenti dell'intera materia. Questa modalità di ripasso in gruppo favorisce una copertura migliore del programma del corso rispetto ad uno studio individuale, nel quale alcuni argomenti potrebbero essere involontariamente trascurati.
Lo scopo di questo progetto è quindi quello di cercare di ricreare le dinamiche appena indicate nel caso in cui non sia possibile un incontro di gruppo tra gli studenti.


\section{Stato dell'Arte}
 Allo stato attuale, sono molte le varianti di giochi a quiz disponibili sul mercato. La motivazione principale risiede nel fatto che di questa tipologia di gioco ne è stato fatto largo uso fin dall'antichità \cite{quizgame} e, successivamente, con l'avvento della televisione ha avuto una seconda ondata di interesse \cite{gameshow}. Le varianti più di rilievo mantengono sempre il concetto di rispondere ad una domanda o scegliendo tra più risposte quella corretta \cite{whowantstobeamillionaire}, o con degli aiuti di diversa natura, o con un tempo limite \cite{jeopardy}. 
 Anche dal punto di vista didattico, sono state sviluppate piattaforme, ad esempio \href{https://quizlet.com/}{Quizlet}, ma le soluzioni trovate sono orientate alla creazione dei quiz da parte degli insegnanti per gli studenti.
 
 L'idea di questo progetto si differenzia, in parte, per: 
\begin{itemize}
    \item la natura accademica orientata a corsi specifici di un determinato corso di laurea e laurea magistrale;
    \item l'utente target, in quanto ne usufruisce direttamente lo studente e non il docente.
\end{itemize}
 
L'intenzione con la quale si vuole portare avanti il progetto è quella di produrre un prodotto superiore a quelli disponibili nel mercato per una specifica necessità riscontrata, non per apportare innovazione tecnologica.
